\documentclass{beamer}
  \usepackage[utf8]{inputenc}
  \usepackage{fancyvrb}
  \usetheme{Montpellier}


  \title{XML}
  \author{Clément Roque et Bastien Tramoni}\institute{Université des Sciences Montpellier 2, licence d'informatique}
  \date{11 octobre 2012}
  \logo{\includegraphics[height=5mm]{logo.png}}

  \begin{document}



 \begin{frame}

  \titlepage

  \end{frame}
\section*{Sommaire}
\begin{frame}
	\frametitle{Vous avez dit XML ?}	
	\begin{itemize}
 \pause\item Qu'est ce que le XML ?
 \pause\item A quoi est ce que ça ressemble ?
 \pause\item A quoi ça peut bien servir ?
 \pause\item Comment c'est utilisé ?
	\end{itemize}
\end{frame}

\section{Qu'est ce que le XML ?}
\begin{frame}
      \frametitle{Qu'est ce que le XML ?}
    \begin{block}{ Définition }
	\begin{itemize}
  \pause \item Extensible Markup Languages (Langage de Balisage Extensible)
  	\end{itemize}
	\end{block}
	\pause c'est un dérive du SGML (Standard Generalized Markup Language) (1970)
 \end{frame}

\begin{frame}
     
    \begin{block}{Histoire}
	\begin{itemize}
   \pause\item XML Working Group (1996)
   \pause\item W3C (World Wide Web Consortium)
   \pause\item XML devient un meta-langage reconnu a partir de 1998
\end{itemize}

\end{block}
\pause Le but principale d'XML est de rendre SGML plus simple d'utilisation et d'interpretation
\end{frame}


\begin{frame}
    \begin{block}{XML est différent de HTML}
	\begin{itemize}
   \pause\item XML = structuration des données
   \pause\item HTML = representation des données \pause
\end{itemize}

\end{block}

\begin{block}{XML : Stockage d'information(s)}
	\begin{itemize}
	\pause\item Organiser une information, une donnée 
	\pause\item Syntaxe est très simple 
	\pause\item Utilisation nombreuse
\end{itemize}
\end{block}
\end{frame}

\section{A quoi ça ressemble ?}
\begin{frame}

\frametitle{A quoi ça ressemble ?}

	\begin{block}{Syntaxe}
		\begin{enumerate}
		\pause\item Prologue
		\pause\item DTD (Déclaration de Type de Document) \pause (facultatif)
		\pause\item L'arbre des elements
		\end{enumerate}

	\end{block}
\end{frame}


\begin{frame}[containsverbatim]


	\begin{block}{Exemple de syntaxe}
		\begin{Verbatim}[fontsize=\scriptsize]
<?xml version="1.0" encoding="ISO-8859-1"?>

<Racine>
	<NoeudA attr1="val1">
		<contenu attr="A"/>
	</NoeudA>

	<NoeudB attr1="val1" attr2="val2">
		<contenu attr=B"/>
	</NoeudB>
</Racine>
		\end{Verbatim}
	\end{block}
\end{frame}
\section{A quoi ça peut bien servir ?}
\begin{frame}[containsverbatim]
\frametitle{A quoi ça peut bien servir ?}

	\begin{exampleblock}{Exemple concret}
		\begin{Verbatim}[fontsize=\scriptsize]
<?xml version="1.0" encoding="ISO-8859-1"?>

<listecontacts proprietaire="Wolverine"> 
	<contact nom="Wayne" prenom="Bruce"> 
		<numero tel="0707070707"/>
	</contact> 

	<contact nom="Parker" prenom="Peter"> 
		<numero tel="4545454545"/>
	</contact> 
</listecontacts>
			
		\end{Verbatim}
	\end{exampleblock}
\end{frame}

\section{Comment j'exploite mes données ?}
\begin{frame}
\frametitle{Comment j'exploite mes données ?}
\pause
	
	\begin{block}{Grace aux parseurs !}
	\begin{itemize}
	\pause\item Analyse syntaxique
	\end{itemize}
	\end{block}

	\begin{block}{Et grace au DOM !}
	\begin{itemize}
	\pause\item Document Object Model
	\pause\item Construction d'un arbre en mémoire a partir d'un document XML
	\pause\item Methodes d'accès et de modification de cet l'arbre
	\end{itemize}
	\end{block}
	

\end{frame}

\begin{frame}
	
	\begin{block}{Langages de gestion de données}
	\begin{itemize}

	\pause\item Gerer base de données (Xquery)
	\pause\item Des API (SAX)
	\pause\item Web avec (messagerie instantanée, formulaire, animation 2D,...)

	\end{itemize}
	\end{block}
	

\end{frame}



\section{Pour résumer}
\begin{frame}
\frametitle{Pour résumer}
	\pause
	\begin{block}{Un excellent format d'échange}
	\pause Création dune structure descriptive adaptée aux données à traiter
	\end{block}

	\pause
	\begin{block}{Langage simple}
	\pause Proche de la syntaxe de HTML
	\end{block}

	\pause
	\begin{block}{Membre d'une grande famille de technologies}
	\pause Utilisation de ces données pour des taches particulières et choisi par le programeur
	\end{block}

	
	

\end{frame}

\section{Nos sources}
\begin{frame}
\frametitle{Nos Sources}
	\begin{itemize}
	\item www.w3c.org
	\item www.devellopez.com
	\item xml.chez.com/initiation
	\end{itemize}

	
	

\end{frame}







 \end{document} 
